% Standalone problem document, generated from the adjacent README.md.
% Compile with XeLaTeX. Mathematical content is maintained in Markdown.
\documentclass[11pt,a4paper]{article}
\usepackage[fontsize=10.5pt]{scrextend}
\usepackage[margin=22mm,headheight=15pt,headsep=9mm,footskip=11mm]{geometry}
\usepackage{fontspec}
\setmainfont{texgyrepagella}[Extension=.otf,UprightFont=*-regular,BoldFont=*-bold,ItalicFont=*-italic,BoldItalicFont=*-bolditalic]
\setsansfont{texgyreheros}[Extension=.otf,UprightFont=*-regular,BoldFont=*-bold,ItalicFont=*-italic,BoldItalicFont=*-bolditalic]
\setmonofont{texgyrecursor}[Extension=.otf,UprightFont=*-regular,BoldFont=*-bold,ItalicFont=*-italic,BoldItalicFont=*-bolditalic,Scale=MatchLowercase]
\usepackage{amsmath,amssymb}
\usepackage{unicode-math}
\setmathfont{texgyrepagella-math.otf}
\usepackage{xcolor,microtype,enumitem,needspace,fancyhdr}
\usepackage{bookmark}
\definecolor{ink}{HTML}{17344B}
\definecolor{accent}{HTML}{197D80}
\definecolor{muted}{HTML}{596773}
\hypersetup{colorlinks=true,linkcolor=accent,urlcolor=accent,pdftitle={MF-09 - Decidability
of strict stability for rational matrix
families},pdfauthor={Open Problems in Numerical Linear Algebra}}
\urlstyle{same}
\setlength{\parindent}{0pt}
\setlength{\parskip}{5pt plus 1pt minus 1pt}
\linespread{1.04}
\setlength{\emergencystretch}{3em}
\widowpenalty=10000
\clubpenalty=10000
\displaywidowpenalty=10000
\allowdisplaybreaks[1]
\setlist{nosep,leftmargin=1.5em}
\providecommand{\tightlist}{\setlength{\itemsep}{0pt}\setlength{\parskip}{0pt}}
\providecommand{\pandocbounded}[1]{#1}
\pagestyle{fancy}
\fancyhf{}
\fancyhead[L]{\sffamily\footnotesize\color{muted}OPEN PROBLEMS IN NUMERICAL LINEAR ALGEBRA}
\fancyhead[R]{\sffamily\bfseries\footnotesize\color{accent}MF-09}
\fancyfoot[L]{\parbox[b]{0.9\textwidth}{\sffamily\fontsize{8}{10}\selectfont\color{muted}Editorial ratings; a bounded search cannot certify that a problem remains unsolved.\\Literature check: 2026-09-10}}
\fancyfoot[R]{\sffamily\footnotesize\color{muted}\thepage}
\renewcommand{\headrulewidth}{0.3pt}
\renewcommand{\headrule}{\hbox to\headwidth{\color{accent}\leaders\hrule height \headrulewidth\hfill}}
\makeatletter
\renewcommand{\section}{\Needspace{5\baselineskip}\@startsection{section}{1}{0pt}{12pt plus 2pt minus 2pt}{4pt}{\sffamily\bfseries\large\color{ink}}}
\renewcommand{\subsection}{\Needspace{4\baselineskip}\@startsection{subsection}{2}{0pt}{9pt plus 1pt minus 1pt}{3pt}{\sffamily\bfseries\normalsize\color{ink}}}
\makeatother
\setcounter{secnumdepth}{0}
\begin{document}
{\sffamily\small\color{muted}Matrix functions and stability\par}
\vspace{3pt}
{\raggedright\hyphenpenalty=10000\exhyphenpenalty=10000\sffamily\bfseries\fontsize{21}{24}\selectfont\color{ink}Decidability
of strict stability for rational matrix families\par}
\vspace{7pt}
{\sffamily\small\textbf{Difficulty:} extreme\quad\textbf{Importance:} broadly
interesting\par}
{\sffamily\small\bfseries\color{accent}Status: Open\par}
\vspace{4pt}
{\color{accent}\hrule height 0.8pt}
\vspace{6pt}
\textbf{Rating rationale:} Extreme because termination at an unpromised
strict threshold is a foundational decidability barrier; broad impact
includes switched control and algorithmic matrix-product theory.

\subsection{Context and notation}\label{context-and-notation}

The joint spectral radius of a nonempty compact set
\(\mathcal M\subset\mathbb C^{d\times d}\) is

\[
\widehat\rho(\mathcal M)=\lim_{k\to\infty}
\max_{A_1,\ldots,A_k\in\mathcal M}\|A_k\cdots A_1\|_2^{1/k}.
\]

This definition also applies to finite real matrix sets. The ordinary
spectral radius of one matrix is written \(\rho(A)\).

\subsection{Problem statement}\label{problem-statement}

Does a Turing machine exist which, on every input consisting of a finite
nonempty list of matrices in \(\mathbb Q^{d\times d}\), encoded by
binary integer numerators and positive denominators, halts and correctly
decides whether

\[
\widehat\rho(\mathcal M)<1?
\]

Dimension and list length are part of the input. No separation from the
threshold is promised. The question asks for termination and
correctness, without a polynomial running-time requirement.

\subsection{References and status
evidence}\label{references-and-status-evidence}

Jungers,
\href{https://perso.uclouvain.be/raphael.jungers/sites/default/files/kcfinder/files/book.pdf}{\emph{The
Joint Spectral Radius: Theory and Applications}}, §2.2.3, Open Question
1, printed p.~29 of the author manuscript. Blondel and Tsitsiklis,
\href{https://www.sciencedirect.com/science/article/abs/pii/S0167691100000499}{The
boundedness of all products of a pair of matrices is undecidable},
Systems \& Control Letters 41 (2000), 135--140, §2, proves
undecidability of the non-strict test \(\widehat\rho\le1\). That theorem
does not answer the strict test.

The screen included ``joint spectral radius'' with ``strict'',
``decidability'', ``less than one'', and 2025/2026. No deciding
algorithm or applicable undecidability reduction was located. Some later
literature summarizes the non-strict theorem as a generic stability
impossibility; admission here follows its actual threshold.
\href{https://github.com/ajt60gaibb/OpenProblemsInNLA/blob/main/matrix-functions-and-stability/MF-04/README.md}{MF-04}
is a proposed finite-product property, whereas this entry asks for a
decision procedure even if such a property fails.

\subsection{Audit --- 2026-09-10}\label{audit-2026-09-10}

Rechecked the indexed text of
\href{https://perso.uclouvain.be/raphael.jungers/sites/default/files/kcfinder/files/book.pdf}{Jungers,
Open Question 1} and strict-stability/decidability searches. No deciding
algorithm or applicable undecidability theorem was located.
\href{https://websites.umich.edu/~lagarias/doc/corr-infprod.pdf}{Daubechies--Lagarias'
correction} independently distinguishes the unresolved subunit test from
the undecidable unit-threshold problem. Evidence remains historical and
bounded.
\end{document}
